\documentclass[11pt,a4paper]{article}

\usepackage[T1]{fontenc}
\usepackage[utf8]{inputenc}
\usepackage[french]{babel}
\usepackage{lmodern}
\usepackage{geometry}
\usepackage{amsmath}
\usepackage{booktabs}
\usepackage{enumitem}

\geometry{margin=2.5cm}

\title{Vers un moteur d'inf\'erence modulaire pour mod\`eles de langage}
\author{}
\date{}

\begin{document}

\maketitle

\section{Design}

\subsection{Probl\'ematique}

Les moteurs d'inf\'erence actuels, tels que \textit{llama.cpp}, offrent des performances robustes mais reposent majoritairement sur des organisations logicielles fortement coupl\'ees. Cette approche favorise l'efficacit\'e dans un cadre d'usage bien d\'efini, au prix d'une flexibilit\'e limit\'ee d\`es qu'il s'agit d'introduire de nouveaux sch\'emas de calcul, de nouvelles cibles mat\'erielles, ou des strat\'egies de distribution sp\'ecifiques.

Dans ce contexte, notre probl\`eme de recherche peut \^etre formul\'e comme suit : \textit{comment concevoir un moteur d'inf\'erence qui conserve des performances comp\'etitives tout en permettant une composition explicite et extensible des composants de la cha\^ine d'ex\'ecution ?}

\subsection{Limites des approches existantes}

L'analyse des solutions existantes met en \`evidence quatre limitations principales.

\begin{enumerate}[leftmargin=*]
	\item \textbf{Couplage structurel \`a la pipeline d'inf\'erence.} Les \`etapes de chargement, d'optimisation et d'ex\'ecution sont souvent imbriqu\'ees, ce qui complique la substitution d'un composant isol\'e.
	\item \textbf{Architecture impos\'ee.} Les outils contraignent l'utilisateur \`a adopter un sch\'ema d'ex\'ecution pr\'ed\'efini, limitant l'exploration de variantes algorithmiques ou syst\`emes.
	\item \textbf{Co\^ut d'extension \'elev\'e.} L'ajout de nouveaux op\'erateurs, de nouvelles politiques de partitionnement ou de nouveaux backends implique fr\'equemment des modifications transversales du code.
	\item \textbf{Faible testabilit\'e modulaire.} L'\'evaluation d'un composant isol\'e est difficile lorsque l'architecture impose l'ex\'ecution de la pipeline compl\`ete.
\end{enumerate}

\subsection{Hypoth\`ese de conception}

Nous posons l'hypoth\`ese qu'une repr\'esentation interm\'ediaire explicite du calcul, sous forme de graphe orient\'e acyclique (DAG), constitue le bon niveau d'abstraction pour d\'ecoupler les responsabilit\'es du moteur.

Cette hypoth\`ese conduit \`a trois choix structurants :

\begin{itemize}[leftmargin=*]
	\item \textbf{Abstraction du calcul.} Le mod\`ele est projet\'e vers une repr\'esentation DAG ind\'ependante du format source.
	\item \textbf{S\'eparation du ``quoi'' et du ``o\`u''.} Le graphe encode les d\'ependances de calcul, tandis que les modules backend d\'efinissent l'ex\'ecution sur une cible mat\'erielle donn\'ee.
	\item \textbf{Composabilit\'e des transformations.} Les politiques de partitionnement et d'optimisation sont d\'ecrites comme des transformations param\'etrables du graphe.
\end{itemize}

\subsection{Architecture modulaire propos\'ee}

L'architecture est organis\'ee en quatre couches fonctionnelles :

\begin{enumerate}[leftmargin=*]
	\item \textbf{Ingestion mod\`ele} : chargement depuis un format source (e.g., GGUF).
	\item \textbf{IR commune} : conversion vers un DAG unifi\'e (n\oe uds op\'erationnels, ar\^etes de donn\'ees).
	\item \textbf{Optimisation/partitionnement} : application de strat\'egies de d\'ecoupage et d'ordonnancement.
	\item \textbf{Ex\'ecution backend} : projection du plan d'ex\'ecution vers une ou plusieurs cibles (CPU, CUDA, Metal, etc.).
\end{enumerate}

Ce d\'ecouplage permet de faire \'evoluer une couche sans remettre en cause l'ensemble du syst\`eme, et facilite la comparaison contr\^ol\'ee de variantes d'impl\'ementation.

\subsection{Objectifs de conception}

Les objectifs associ\'es \`a cette proposition sont les suivants :

\begin{itemize}[leftmargin=*]
	\item fournir une interface d'inf\'erence extensible et stable ;
	\item permettre le partitionnement explicite du graphe de calcul ;
	\item r\'eduire le co\^ut d'int\'egration de nouveaux op\'erateurs et backends ;
	\item am\'eliorer la testabilit\'e et la reproductibilit\'e exp\'erimentale.
\end{itemize}

\subsection{Synth\`ese}

En r\'esum\'e, la contribution de design ne se limite pas \`a une optimisation locale de l'ex\'ecution, mais propose une structuration logicielle orient\'ee recherche et exp\'erimentation. Le DAG devient l'artefact central \`a partir duquel sont d\'eclin\'ees les d\'ecisions de partitionnement et d'ex\'ecution, ouvrant la voie \`a une \'evaluation comparative plus rigoureuse des strat\'egies d'inf\'erence modulaire.

\section{Implementation}

\subsection{Vue d'ensemble}

L'impl\'ementation est r\'ealis\'ee sous forme de modules C/C++ faiblement coupl\'es, int\'egr\'es \`a GGML. Le choix d'une granularit\'e modulaire vise deux objectifs pratiques : (i) permettre la validation unitaire de chaque composant et (ii) faciliter l'\'evolution du moteur sans r\'e\'ecriture globale.

Les principaux modules sont :
\begin{itemize}[leftmargin=*]
	\item \texttt{parser.c/.h} : chargement d'un mod\`ele GGUF et extraction des m\'etadonn\'ees ;
	\item \texttt{distrib.c/.h} : repr\'esentation du graphe de partitionnement et heuristiques d'\'equilibrage ;
	\item \texttt{dummy\_distrib.c} : baseline de partitionnement ;
	\item \texttt{KaHIP.c} : interface de partitionnement avanc\'e via KaHIP ;
	\item \texttt{parser\_main.c} : point d'entr\'ee exp\'erimental.
\end{itemize}

\subsection{Composant de parsing de mod\`ele}

Le composant \texttt{parser} impl\'emente le chargement GGUF et instancie une structure de mod\`ele contenant le contexte GGML, le contexte GGUF, et les hyperparam\`etres utiles \`a l'inf\'erence.

Fonctions d'int\'er\^et :
\begin{itemize}[leftmargin=*]
	\item \texttt{model\_load(...)} : validation du fichier et initialisation du mod\`ele ;
	\item \texttt{run\_single\_token\_inference(...)} : ex\'ecution minimale pour v\'erification fonctionnelle ;
	\item \texttt{run\_batch\_inference(...)} : ex\'ecution multi-tokens ;
	\item \texttt{save\_logits(...)} : export des sorties pour comparaison.
\end{itemize}

Le format GGUF est exploit\'e comme format canonique d'entr\'ee. En pratique, le parseur suppose la disponibilit\'e d'un en-t\^ete coh\'erent (version, nombre de tenseurs, m\'etadonn\'ees) et de descripteurs de tenseurs correctement nomm\'es et typ\'es.

\subsection{D\'etails du parseur GGUF}

Le parseur est con\c{c}u pour (i) valider la structure du fichier, (ii) construire les tenseurs GGML correspondants et (iii) produire une repr\'esentation op\'erationnelle (DAG) exploitable par l'outil de partitionnement.

Étapes principales :
\begin{enumerate}[leftmargin=*]
	\item Lire l'en-t\^ete GGUF et v\'erifier la magie/version.
	\item Pour chaque descripteur de tenseur, allouer un \texttt{ggml\_tensor} et enregistrer son nom, sa forme et son type.
	\item Construire la table de correspondance \textit{name -> tensor}\ pour retrouver un tenseur par ses noms alternatifs (ex: \texttt{"weight"} / \texttt{"W"}).
	\item Construire le DAG de forward en parcourant la d\'ependance des op\'erateurs dans le contexte GGML (fonctions d'extraction d'alias et d'usage).
\end{enumerate}

Fonctions et points d'impl\'ementation notables :
\begin{itemize}[leftmargin=*]
	\item \texttt{find\_tensor\_any(...)} : renvoie le pointeur vers un tenseur existant parmi plusieurs noms candidats ; utile pour la robustesse vis-\`a-vis de variantes de dump.
	\item \texttt{load\_required\_tensor\_any(...)} : m\'ethode d'abstraction pour v\'erifier et charger les tenseurs indispensables.
\end{itemize}

Exemple de heuristique de mappage : pour estimer la taille m\'emoire d'un tenseur on calcule
\[
\mathrm{bytes}(T) = \prod_{d\in \mathrm{dims}} d \times \mathrm{sizeof}(\text{type})
\]
et on stocke cette valeur comme attribut utilis\'e plus tard pour l'estimation des co\^uts de communication.

\subsection{Construction du graphe de partitionnement}

Le graphe de partitionnement est construits\'a partir du DAG d'ex\'ecution : chaque op\'eration devient un sommet avec un co\^ut de calcul estim\'e, et une ar\^ete est cr\'ee entre deux sommets si un tenseur produit par l'un est consomm\'e par l'autre.

Calcul des co\^uts de sommet (co\^ut de calcul) : une estimation simple utilis\'ee dans l'impl\'ementation est
\[
C_{\mathrm{comp}}(v) = \frac{\mathrm{FLOPs}(v)}{\mathrm{peak\_flops}_{\mathrm{target}}},
\]
o\`u $\mathrm{FLOPs}(v)$ est une estimation du nombre d'op\'erations flottantes n\'ecessaires pour ex\'ecuter $v$ (par exemple, $2\times n \times m$ pour une multiplication matrice-vecteur). Le param\`etre $\mathrm{peak\_flops}_{\mathrm{target}}$ est fourni par l'utilisateur ou obtenu par profilage.

Calcul des poids d'ar\^ete (co\^ut de communication) : pour une ar\^ete transportant un tenseur $T$ de taille $B$ octets,
\[
W_{\mathrm{comm}}(e) = \alpha + \frac{B}{\beta},
\]
o\`u $\alpha$ est la latence fixe (overhead) et $\beta$ la bande passante effective. Ces constantes peuvent \^etre estim\'ees \`a partir de micro-benchmarks.

Pseudocode de construction :
\begin{verbatim}
for each op v in DAG:
	comp_cost[v] = estimate_flops(v) / peak_flops
	for each output tensor t of v:
		for each consumer u of t:
			add edge (v -> u) with weight = alpha + bytes(t)/bandwidth
\end{verbatim}

\subsection{Algorithmes de partitionnement}

\paragraph{Baseline greedy (\texttt{dummy\_distrib})}

La baseline impl\'ement\'ee est une heuristique de bin-packing gloutonne : on ordonne les sommets par co\^ut decroissant et on les assigne \`a la partition qui minimise l'augmentation de l'objectif $\mathcal{L}$. Cette strat\'egie est simple mais robuste pour les petits graphes.

\paragraph{Partitionnement via KaHIP}

Pour KaHIP, le graphe est converti en un format que KaHIP accepte (hypergraph/graph selon l'API), puis l'appel est d\'el\'egu\'e soit via l'API C, soit en invoquant l'ex\'ecutable. Le r\'esultat est une table \texttt{vertex\_to\_partition} qu'on r\'eimporte et utilise pour annoter le plan d'ex\'ecution.

Fichier d'\'echange typique (graph): liste d'ar\^etes avec poids et nombre de sommets, suivie des poids de sommet.

\subsection{Interface avec KaHIP (d\'etails)}

L'int\'egration comporte deux variantes : (i) appel direct de la librairie KaHIP lors de la compilation, (ii) export du graphe sur disque et invocation du binaire KaHIP.

Dans la voie ``API'', on construit les structures C attendues par KaHIP et on appelle \texttt{kaffpa\_interface(...)} (ou l'API appropri\'ee selon la version). Dans la voie binaire, on produit un fichier texte \texttt{graph.kt} et on ex\'ecute :
\begin{verbatim}
kahip --input=graph.kt --k=4 --preconfiguration=fast
\end{verbatim}

Ensuite, le fichier de sortie associant chaque sommet \`a une partition est pars\'e et copi\'e dans \texttt{vertex\_to\_partition}.

\subsection{Tests, benchmarks et protocole de validation}

Pour garantir la reproductibilit\'e, nous proposons un ensemble de tests unitaires et de micro-benchmarks :

\begin{itemize}[leftmargin=*]
	\item \textbf{Unitaires parser} : v\'erifier que le nombre de tenseurs lu correspond \`a la m\'etadonn\'ee et que les noms principaux existent (assertions sur \texttt{n\_tensors}, \texttt{hparams}).
	\item \textbf{Tests de r\'egression logits} : ex\'ecuter \texttt{run\_single\_token\_inference} sur un contexte standard et comparer les logits aux valeurs de r\'ef\'erence (tol\'erance relative, p.ex. $10^{-5}$).
	\item \textbf{Benchmarks} : \texttt{test\_perf} fournit le temps d'ex\'ecution simul\'e par sommet (en ms), la coupe inter-partitions estim\'ee et le ratio d'\'equilibre.
\end{itemize}

Commandes minimales :
\begin{verbatim}
mkdir -p build && cd build
cmake .. -DGGML_USE_KAHIP=ON -DKAHIP_ROOT=/path/to/kahip-install
cmake --build . --target parser test_perf model_compare
// Run parser
./build/parser sample_model.gguf
// Run perf
./build/test_perf sample_model.gguf
\end{verbatim}

Crit\`eres d'acceptation :
\begin{itemize}[leftmargin=*]
	\item charge-structure correcte (DAG sans cycles) ;
	\item diff de logits entre r\'ef\'erence et impl\'ementation $<$ tol ;
	\item partitions valides (aucun sommet non assign\'e), ratio d'\'equilibre raisonnable (p.ex. $<1.2$).
\end{itemize}

\subsection{Notes d'impl\'ementation et limites}

Quelques choix d'impl\'ementation m\'eritent d'\^etre soulign\'es :
\begin{itemize}[leftmargin=*]
	\item l'estimation de FLOPs est approximative et d\'epend du type d'op\'eration ; une instrumentation runtime am\'eliore la pr\'ecision ;
	\item la liaison avec KaHIP suppose une conversion correcte du graphe en structure accept\'ee ; cela peut impliquer une perte d'information sur les co\^uts d'ar\^etes si l'on utilise une repr\'esentation simplifi\'ee ;
	\item l'impl\'ementation actuelle ne g\`ere pas les op\'erateurs dynamiques ni les boucles explicites ; ces cas n\'ecessitent une transformation pr\'ealable du graphe.
\end{itemize}
 
\end{document}
\subsection{Composant de partitionnement}

Le composant \texttt{distrib} transforme le graphe de calcul en graphe pond\'er\'e de partitionnement : les sommets repr\'esentent le co\^ut de calcul des op\'erations, et les ar\^etes repr\'esentent le co\^ut de communication entre op\'erations d\'ependantes.

L'objectif optimis\'e est d\'efini par :
\[
\mathcal{L} = \lambda_{\mathrm{comm}}\,C_{\mathrm{cut}} + \lambda_{\mathrm{bal}}\,P_{\mathrm{balance}},
\]
o\`u $C_{\mathrm{cut}}$ est le co\^ut de coupe inter-partitions et $P_{\mathrm{balance}}$ la p\'enalit\'e de d\'es\'equilibre de charge. Cette formulation permet d'ajuster explicitement le compromis communication/\'equilibrage.

La configuration de partitionnement inclut notamment : le nombre de partitions, les poids de l'objectif, un plafond de charge \textit{hard cap}, et des param\`etres d'it\'eration pour le raffinage local.

\subsection{Int\'egration KaHIP}

Le module \texttt{KaHIP.c} encapsule l'utilisation de KaHIP comme solveur externe de partitionnement. L'int\'egration est optionnelle au moment de la configuration CMake, ce qui permet de comparer deux r\'egimes :
\begin{itemize}[leftmargin=*]
	\item \textbf{baseline interne} (sans KaHIP) ;
	\item \textbf{partitionnement avanc\'e} (avec KaHIP).
\end{itemize}

Cette dualit\'e constitue une base exp\'erimentale utile pour quantifier l'apport d'un partitionneur externe sur des graphes de tailles diff\'erentes.

\subsection{D\'etails de reproductibilit\'e}

\subsubsection{Protocole de validation}

Pour reproduire les r\'esultats techniques de base, nous recommandons la s\'equence suivante :
\begin{enumerate}[leftmargin=*]
	\item v\'erifier le chargement du mod\`ele GGUF via \texttt{parser} ;
	\item contr\^oler la coh\'erence des sorties (logits) sur un cas simple ;
	\item ex\'ecuter les benchmarks (\texttt{test\_perf}) ;
	\item comparer les strat\'egies de partitionnement (\texttt{model\_compare}).
\end{enumerate}

\subsection{Discussion d'impl\'ementation}

L'impl\'ementation actuelle constitue un socle reproductible pour l'\'etude du partitionnement de graphes d'inf\'erence, mais ne couvre pas encore l'ex\'ecution distribu\'ee compl\`ete multi-backends. En cons\'equence, les r\'esultats obtenus \`a ce stade doivent \^etre interpr\'et\'es comme des r\'esultats de faisabilit\'e et de caract\'erisation, plut\^ot que comme une validation syst\`eme de bout en bout.

\end{document}
